\chapter{Simplicity on Elements: The Unoccupied Optimal Corner}
\label{app:simplicity}

\section{Scope and Motivation}
\label{app:simplicity:scope}

Simplicity~\cite{OC17} is an \emph{alternative script primitive} --- a
typed combinator calculus with Coq-verified denotational semantics ---
deployed on Elements~\cite{DPNS16}, a Bitcoin-adjacent sidechain whose
blocks are produced by known functionaries rather than by proof-of-work
miners. Execution uses native-code \emph{jets} for known combinator
subtrees. Because Elements' consensus model differs from Bitcoin's,
we exclude Simplicity from the formal results of
Chapter~\ref{ch:security} and report it here as a cross-substrate
reference: Simplicity demonstrates that the dominant corner $D^\star$
of Chapter~\ref{ch:security} is occupiable on some substrate, even
while no Bitcoin covenant proposal currently occupies it.

This distinction is load-bearing for the thesis's claim.
Chapter~\ref{ch:security}'s contribution is \emph{unoccupied on
Bitcoin}, not \emph{theoretically unoccupiable}. Simplicity
is the constructive existence proof for the latter. A reader who
conflates the two would mistake a specification gap for an
impossibility theorem.

\section{Design-Space Position}
\label{app:simplicity:axes}

Our Simplicity vault occupies the strict-best axis values on every
axis where such a value exists, in the vocabulary of
Section~\ref{se:design_space_axes}:

\begin{itemize}
  \item \axauth{} = \texttt{cold-key}: key-gated recovery via a
    Schnorr signature check.
  \item \axrevault{} = \texttt{no}: atomic withdrawal; the program
    enforces the complete output tuple via \texttt{jet::outputs\_hash()
    = param::VAULT\_RECOVER\_OUTPUTS\_HASH} (see
    \texttt{simf/vault.simf}).
  \item \axfee{} = \texttt{dynamic} (Elements fee structure): fees
    live inside the transaction; no out-of-band anchor.
  \item \axwdbind{} = \texttt{dedicated-jets-combinator}: the
    withdrawal destination is bound through a typed output-hash
    commitment at compile time, stronger than any of the
    Bitcoin proposals on the same axis.
  \item \axrecbind{} = \texttt{dedicated-jets-combinator}: same
    binding for the recovery leaf, eliminating class \classcold{}.
  \item \axopsurf{} = \texttt{single-mode}: a program's meaning is
    determined by its Merkle root; no \texttt{OP\_SUCCESSx}-style
    fallthrough surface.
  \item \axintro{} = \texttt{typed-combinator-introspection}: the
    fundamental operators are total, typed, and Coq-verified.
\end{itemize}

Consequently, Simplicity's vault program is immune by construction
to classes \classfee{}, \classgrief{}, \classscale{}, \classhot{},
\classcold{}, and \classmode{} --- every attack class identified in
Chapter~\ref{ch:security}. It occupies the strict-best position on
\axauth{}, \axrevault{}, \axwdbind{}, \axrecbind{}, and \axintro{}
simultaneously; no Bitcoin proposal does.

\section{Lifecycle Measurements}
\label{app:simplicity:empirical}

The Simplicity vault lifecycle totals $865\vB{}$ (deposit $269$,
trigger $298$, withdraw $298$, recover $286$), $49$--$135\%$ larger
than the Bitcoin designs of Chapter~\ref{ch:results}. The overhead
bundles Elements' serialisation (${\sim}30$--$40$ bytes per output
for confidential-transaction fields even in non-confidential mode)
and Simplicity's bit-oriented DAG witness format. Elements
functionaries do not operate a public mempool, so class \classfee{}
(fee-channel pinning) has no substrate analogue under
honest-federation assumptions.

Our Simplicity vault is atomic: recovery binds the complete output
tuple via the jets commitment cited above, admitting neither partial
withdrawal nor batched recovery and placing the program at
\axrevault{} = \texttt{no} in the design-space lattice. Class
\classscale{} (per-UTXO recovery-cost scaling) has no surface.
Cold-key compromise (TM11) forces recovery but cannot redirect
funds, matching CTV and OP\_VAULT. A different Simplicity program
could express \axrevault{} = \texttt{yes} by replacing the
whole-tuple commitment with per-output checks; that design would
inherit class \classscale{} under the D2 implication
(Section~\ref{se:design_space_axes}) but is not what we measured.

\section{Foundations and Substrate}
\label{app:simplicity:foundations}

Simplicity's Coq-verified denotational semantics give stronger
guarantees than Alloy bounded model checking of the Bitcoin Script
extensions can produce --- a Simplicity program's meaning is
determined by its Merkle root, under a formal semantics that has
been mechanically verified. However, Simplicity on Liquid is live
today under \emph{federation trust}, whereas the four Bitcoin Script
extensions remain under \emph{Bitcoin consensus review}.

The substrate choice is a security-relevant decision alongside the
per-covenant analysis of Chapter~\ref{ch:security}: a vault on
Simplicity inherits the federation's liveness and censorship
properties; a vault on Bitcoin inherits proof-of-work's. The
dominant corner being occupiable on Simplicity is therefore not a
turn-key prescription to migrate vaults to Liquid; it is evidence
that the unoccupied status among current Bitcoin proposals reflects a
consensus-substrate constraint rather than a design impossibility.

\section{Future Work: Porting the Optimal Corner to Bitcoin}
\label{app:simplicity:future}

A Simplicity-style vault on Bitcoin would require Simplicity's
taproot leaf type, which is specified in a draft BIP but not under
active Bitcoin consensus review. Closing the Simplicity-to-Bitcoin
gap depends on that soft-fork path. An alternative that preserves
Bitcoin's consensus substrate is a CAT+CSFS-based bound-recovery
construction (Poelstra's sketch~\cite{Poe21}) that moves CAT+CSFS
to \axrecbind{} = \texttt{signature-dual-bind}; this would eliminate
class \classcold{} while retaining the hot-key-theft immunity of
class \classhot{}, placing CAT+CSFS at the dominant corner of the
Bitcoin sublattice. Neither path is available as of
thesis submission; we leave both as open problems.
